\section{Compiling C Data Structures}

\subsection{Basic Data Types and Machine Representation}

\begin{definition}
    \textbf{Primitive Data Types in x86-64 Linux (System V ABI):} Primitive types have fixed sizes and alignment requirements:
    \begin{center}
        \renewcommand{\arraystretch}{1.15}
        \setlength{\tabcolsep}{6pt}
        \begin{tabular}{|l|c|c|c|l|}
            \hline
            \textbf{Type} & \textbf{Size} & \textbf{Alignment} & \textbf{Suffix} & \textbf{Description} \\
            \hline
            \texttt{char} & 1 B & 1 B & \texttt{b} & Byte \\
            \texttt{short} & 2 B & 2 B & \texttt{w} & Word (16 bits) \\
            \texttt{int} & 4 B & 4 B & \texttt{l} & Doubleword (32 bits) \\
            \texttt{long} & 8 B & 8 B & \texttt{q} & Quadword (64 bits) \\
            \texttt{long long} & 8 B & 8 B & \texttt{q} & Quadword (64 bits) \\
            \texttt{T *} (pointer) & 8 B & 8 B & \texttt{q} & Quadword address \\
            \texttt{float} & 4 B & 4 B & \texttt{s} / \texttt{ss} & Single precision (SSE) \\
            \texttt{double} & 8 B & 8 B & \texttt{l} / \texttt{sd} & Double precision (SSE) \\
            \texttt{long double} & 16 B & 16 B & \texttt{t} & Extended precision \\
            \hline
        \end{tabular}
    \end{center}
\end{definition}

\begin{important}
    Assembly code does not retain high-level C types. All data is viewed merely as untyped byte sequences in registers or flat virtual memory.
\end{important}

\subsection{One-Dimensional Arrays}

\begin{definition}
    An array declaration \texttt{T A[N]} allocates a \textbf{contiguous} block of $N \times L$ bytes, where $L = \text{sizeof}(T)$. The identifier \texttt{A} acts as a pointer pointing to the first element (\texttt{\&A[0]}).
\end{definition}

\begin{important}
    C does not perform array bounds checking at compile-time or runtime. Accessing out-of-bounds indices reads or corrupts adjacent memory or triggers segmentation faults.
\end{important}

\begin{definition}
    \textbf{Indexing} \texttt{A[i]} is semantically identical to pointer dereferencing \texttt{*(A + i)}. Given starting address $x_A$, the memory address of element $i$ is $\&A[i] = x_A + L \cdot i$
\end{definition}

\begin{remark}
    \textbf{Assembly Translation:} Memory operands in x86-64 map directly to array indexing using the scaled index addressing mode $D(R_b, R_i, S)$:
    $$\text{Mem}[\text{Reg}[R_b] + S \cdot \text{Reg}[R_i] + D]$$
    where $S \in \{1, 2, 4, 8\}$ corresponds to element size $L$. The address of an element can be computed without reading memory using \texttt{leaq}.
\end{remark}


\subsection{Nested Arrays}

\begin{definition}
    \textbf{A multidimensional array} declaration \texttt{T A[R][C]} allocates a single contiguous memory block of $R \times C \times L$ bytes, where $R$ is rows, $C$ is columns, and $L = \text{sizeof}(T)$.
\end{definition}

\begin{remark}
     C stores nested arrays row-by-row sequentially. All elements of row $0$ are stored first, followed contiguously by row $1$, row $2$, and so forth.
\end{remark}

\begin{theorem}
    \textbf{Address Computation:} The starting address of row $i$ (\texttt{A[i]}) is $\&A[i][0] = x_A + i \cdot (C \cdot L)$ and the memory address of element \texttt{A[i][j]} is $\&A[i][j] = x_A + L \cdot (i \cdot C + j) = x_A + (i \cdot C \cdot L) + (j \cdot L)$
\end{theorem}

\begin{remark}
    \textbf{Compiler Optimizations:} When the column count $C$ is a constant or a power of two, the compiler avoids costly multiplication instructions (\texttt{imulq}) by replacing $i \cdot C$ with bit shifts (\texttt{salq}) and scaled addition (\texttt{leaq}).
\end{remark}

\subsection{Multi-Level Arrays}

\begin{definition}
    A \textbf{multi-level array} declaration \texttt{T *B[R]} allocates an array of $R$ contiguous pointers (each 8 bytes on x86-64), where each entry \texttt{B[i]} holds a pointer to a separate, independently allocated 1D array of elements of type \texttt{T}.
\end{definition}

\begin{important}
    Rows are not required to be contiguous to each other and can have variable lengths (ragged/jagged arrays).
\end{important}

\subsection{Structures and Data Alignment}

\begin{definition}
    \textbf{Structures (\texttt{struct}):}
    \begin{itemize}
        \item A \texttt{struct} groups heterogeneous data fields into a single contiguous memory block.
        \item Field order is \textbf{strictly preserved} as declared in the C source code.
        \item The compiler computes fixed byte offsets from the base address for each field at compile time.
        \item Accessing a struct field compiles to base-plus-displacement addressing: \texttt{movl D(\%rdi), \%eax}.
    \end{itemize}
\end{definition}

\begin{remark}
    Modern CPU memory architectures and cache lines (typically 64 bytes) fetch data from aligned addresses (multiples of 2, 4, 8, or 64 bytes). If data crosses a cache line or virtual page boundary, accessing it requires multiple memory bus cycles, severely degrading memory bandwidth and CPU performance.
\end{remark}

\begin{theorem}
    \textbf{K-byte Alignment Rule:} Any primitive object of size $K$ bytes must reside at a memory address that is a multiple of $K$:
    $$\text{Address} \equiv 0 \pmod K$$
\end{theorem}

\begin{theorem}
    \textbf{Structure Alignment and Padding Rules:}
    \begin{itemize}
        \item \textbf{Overall Struct Alignment ($K_{\text{struct}}$):} The alignment of the entire struct is dictated by the maximum alignment required by any of its members:
        $$K_{\text{struct}} = \max_i(K_{\text{field}_i})$$
        The struct's starting address must be a multiple of $K_{\text{struct}}$.
        \item \textbf{Internal Padding:} Unused gap bytes inserted between adjacent fields so that each individual field starts at an offset satisfying its own alignment constraint.
        \item \textbf{External (Trailing) Padding:} Unused gap bytes appended to the end of the struct so that the total struct size is an exact multiple of $K_{\text{struct}}$.
    \end{itemize}
\end{theorem}

\begin{important}
    \textbf{Why External Padding is Mandatory:} External padding ensures that when structs are placed into an array (\texttt{struct S arr[N]}), every element \texttt{arr[i]} starts at an offset that satisfies $K_{\text{struct}}$. Without trailing padding, the second element \texttt{arr[1]} would violate alignment.
\end{important}

\begin{remark}
    Because field order is preserved by the compiler, poor ordering introduces wasted internal padding. Ordering fields in \textbf{descending order of size} eliminates or minimizes internal padding.
\end{remark}

\subsection{Unions}

\begin{definition}
    \textbf{Unions (\texttt{union}):}
    \begin{itemize}
        \item A \texttt{union} allocates all of its member fields at the exact same starting memory address (offset 0), effectively overlaying them.
        \item Only one member field can be validly stored at any given point in time.
        \item \textbf{Size and Alignment:} The total size of a union equals the size of its largest member, rounded up to a multiple of its largest alignment constraint ($K_{\text{union}} = \max_i(K_{\text{field}_i})$).
    \end{itemize}
\end{definition}